\section{The appointed elements}
\begin{properinventory}{Textual unit}{Text / citation}
\textbf{Entrance and prayer} & \\
Entrance & \textit{Ego clamavi}; cf. Psalm 16:6,8 (Vulgate), 17:6,8 (Hebrew); Missal-antiphon route.\\
Collect & \textit{Omnipotens sempiterne Deus, fac nos}; Week XXIX.\\
\textbf{Liturgy of the Word} & \\
First reading & Isaiah 45:1,4--6; Cyrus and the one Lord.\\
Responsorial Psalm & Psalm 96:1,3,4--5,7--8,9--10a,c (Hebrew; Vulgate 95); response 7b. The middle clause of verse 10 is not sung.\\
Second reading & 1 Thessalonians 1:1--5b; stopping before the final clause about the missionaries' conduct.\\
Alleluia verse & Philippians 2:15d,16a; two short verse portions.\\
Gospel & Matthew 22:15--21; one form, no shorter alternative.\\
\textbf{Eucharistic liturgy} & \\
Prayer over the Offerings & \textit{Tribue nos, Domine, quaesumus}; Week XXIX.\\
Communion & Either \textit{Ecce oculi Domini}, cf. Psalm 32:18--19 (Vulgate), 33:18--19 (Hebrew), or \textit{Filius hominis venit}, Mark 10:45, ransom excerpt. Neither option is selected here.\\
Prayer after Communion & \textit{Fac nos, quaesumus, Domine}; Week XXIX.\\
\end{properinventory}

The Psalms at Entrance and Communion are adaptations; the Mark antiphon is an excerpt, not the whole verse. Approved chant choices can replace the Missal-antiphon route. No particular local chant, Preface or Eucharistic Prayer has been selected. The Sunday Gloria and Creed belong to the Ordinary. The precise books, permitted branches and occurrence appear in the liturgical appendix.
